\chapter{Piezoelektrizität und Thermopiezoelektrizität}\label{chap:coupling}
\section{Stoffgesetz, Enthalpie und Variation}
Das lineare Stress-Charge-Gesetz mit optionaler thermischer Eigenverzerrung ist
\begin{equation}\label{eq:piezo-law}
\sigma_v=\mathbf C^E\varepsilon_m-\mathbf e^T\mathbf E,\qquad
\mathbf D=\mathbf e\varepsilon_c+\boldsymbol\kappa^S\mathbf E,
\quad\varepsilon_c=\varepsilon_v-\chi\varepsilon_{\rm th}.
\end{equation}
$\chi=0$ entspricht der Standardwahl \texttt{electric\_thermal\_strain=False};
$\chi=1$ dem ausdrücklich aktivierten Term. Mechanik verwendet immer
$\varepsilon_m=\varepsilon_v-\varepsilon_{\rm th}$.
Ein pyroelektrischer Term $\mathbf p\theta$ wird nicht hinzugefügt.
\index{Piezoelektrizität}\index{Elektrische Enthalpie}

Bei fest bekannter Temperatur ist die implementierte Enthalpiedichte
\begin{equation}\label{eq:enthalpy}
\mathcal h=\tfrac12\varepsilon_m^T\mathbf C^E\varepsilon_m
-\varepsilon_c^T\mathbf e^T\mathbf E
-\tfrac12\mathbf E^T\boldsymbol\kappa^S\mathbf E.
\end{equation}
$\partial\mathcal h/\partial\varepsilon_v=\sigma_v$ und
$-\partial\mathcal h/\partial\mathbf E=\mathbf D$.
Mit $\delta\mathbf E=-\nabla\delta\phi$ folgt
\begin{equation}\label{eq:enthalpy-variation}
\delta\mathcal H=\int_\Omega
\varepsilon_v(\delta\mathbf u)^T\sigma_v+
\nabla\delta\phi^T\mathbf D\,\mathrm dV.
\end{equation}
Die Variationen werden mit mechanischer Arbeit und der signierten elektrischen
rechten Seite aus Gleichung~\ref{eq:electro-weak} gleichgesetzt.
Daraus entstehen gleichzeitig mechanische und elektrische schwache Form.
Der negative elektrische Energieanteil ist die Legendre-transformierte
Enthalpie für Potentialvariablen; er bedeutet weder negative Permittivität
noch eine negative gespeicherte Dielektrikumsenergie.
Diese eigene Herleitung verwendet die Projektvorzeichen; die externe
Theoriequelle \cite{comsol-piezo} dient nur zum Abgleich der Stoffgesetzformen.

\section{Elementblöcke und Symmetrie}
Durch $\varepsilon_v=\mathbf B\mathbf u_e$ und $\mathbf E=-\mathbf G\phi_e$ folgt
\begin{align}\label{eq:piezo-blocks}
\mathbf K_u^e&=\int_K\mathbf B^T\mathbf C^E\mathbf B,&
\mathbf L^e&=\int_K\mathbf B^T\mathbf e^T\mathbf G,\\
\mathbf D^e&=\int_K\mathbf G^T\boldsymbol\kappa^S\mathbf G,&
\mathbf A^e&=\begin{pmatrix}\mathbf K_u^e&\mathbf L^e\\
(\mathbf L^e)^T&-\mathbf D^e\end{pmatrix}.\nonumber
\end{align}
Die Maße sind $\mathrm dV$. Blöcke haben N/m, N/V (= C/m), F.
Tet4 multipliziert die konstanten Integranden mit $V$, Tet10 integriert
punktweise. Gleiche Quadratur für die Kopplungsblöcke erhält Reziprozität.
Ein konstantes Potential hat $\mathbf G\mathbf 1=0$ und erzeugt keine Kraft.
Der Sattelpunktoperator ist symmetrisch-indefinit; CG ist ungeeignet.

Für die thermische rechte Seite gilt
\begin{equation}\label{eq:piezo-thermal-rhs}
\mathbf f_{\rm th}^e=\begin{pmatrix}
\int_K\mathbf B^T\mathbf C^E\varepsilon_{\rm th}\,\mathrm dV\\
\chi\int_K\mathbf G^T\mathbf e\varepsilon_{\rm th}\,\mathrm dV
\end{pmatrix}.
\end{equation}
Der positive elektrische RHS folgt aus dem negativen bekannten D-Term;
er darf nicht isoliert umgedreht werden. Unter einer einzigen gemeinsamen
Referenztemperatur kann man den Ablauf schematisch als
\begin{equation}\label{eq:triangular-coupling}
\begin{pmatrix}
\mathbf K_u&\mathbf K_{uT}&\mathbf L\\
0&\mathbf K_T&0\\
\mathbf L^T&\mathbf K_{\phi T}&-\mathbf D
\end{pmatrix}
\begin{pmatrix}\mathbf u\\\boldsymbol\theta\\\boldsymbol\phi\end{pmatrix}
=\begin{pmatrix}\mathbf f_u\\\mathbf f_T\\\mathbf f_\phi\end{pmatrix},
\end{equation}
mit $\mathbf K_{uT}=-\int\mathbf B^T\mathbf C^E\alpha\mathbf N$ und
$\mathbf K_{\phi T}=-\chi\int\mathbf G^T\mathbf e\alpha\mathbf N$ schreiben.
Bei regionsweisen Referenztemperaturen wird kein globaler $\theta$-Vektor
vorgegeben: Die Produktionspipeline löst absolute $T$ und bildet die bekannten
zellweisen RHS-Terme. Die Nullblöcke in der Wärmezeile erklären
$T\rightarrow(\mathbf u,\phi)$ und die fehlende thermische Rückwirkung.
\index{Thermopiezoelektrizität}

\begin{figure}[htbp]
\centering
\setlength{\unitlength}{1mm}
\begin{picture}(125,32)
\put(0,10){\framebox(28,14){Wärme $T$}}
\put(28,17){\vector(1,0){28}}
\put(32,21){\makebox(20,5){$\alpha\theta$}}
\put(56,10){\framebox(30,14){Mechanik $u$}}
\put(86,20){\vector(1,0){11}}
\put(97,14){\vector(-1,0){11}}
\put(91,25){\makebox(4,5){$e$}}
\put(97,10){\framebox(28,14){Potential $\phi$}}
\put(29,2){\makebox(95,5){Thermischer D-Term nur bei $\chi=1$}}
\end{picture}
\caption{Einseitiger Wärmeablauf und reziproke Piezo-Kopplung.
Eigene schematische Darstellung; kein zusätzlicher thermischer Rückweg.}
\label{fig:coupling-flow}
\end{figure}

\section{Energieberichte und Randbedingungen}
Berichtet werden mechanische Energie, Kopplungsenthalpie und negative
Dielektrikumsenthalpie getrennt. Ihre Summe ist $\mathcal H$ aus
Gleichung~\ref{eq:enthalpy}. Eine indefinite Enthalpie darf nicht als Norm eines
Diskretisierungsfehlers dienen; dafür werden positive entkoppelte
Gradientennormen verwendet. Tet4 benutzt die Zellmittel-Eigendehnung gemäß
Gleichung~\ref{eq:thermal-variance}; Tet10-Energien sind separate Integrale.

Eine leitende feste Elektrode hat \emph{ein} Potential, eine schwebende Elektrode
\emph{eine} vorgegebene Gesamtleiterladung. Punktweise $\mathbf D\cdot n=0$
ist ein isolierter Rand und allgemein kein Leiter-Leerlauf. Die algebraische
Kondensation in Kapitel~\ref{chap:algebra} sorgt für diese Unterscheidung.
Mechanische Lagerung und elektrische Referenz sind pro Komponente notwendig.
Eine reine Ladungsvorgabe erdet nicht implizit.

\section{Handrechnung: affiner Piezo-Patch}
Auf dem Einheits-Tet4 seien $C_{xx}=10$ Pa, $e_{xx}=2$ C/m$^2$,
$\kappa_{xx}=3$ F/m, $\varepsilon_{xx}=0{,}2$ und
$\partial_x\phi=0{,}3$ V/m. Ohne thermische Eigenverzerrung ist $E_x=-0{,}3$
V/m. Aus Gleichung~\ref{eq:piezo-law} folgen
\begin{equation}\label{eq:piezo-patch-example}
\sigma_{xx}=2{,}6\ {\rm Pa},\qquad D_x=-0{,}5\ {\rm C/m^2}.
\end{equation}
Die Energieterme sind $1/30$ J, $1/50$ J und $-9/400$ J;
$\mathcal H=37/1200$ J. Der mechanische Term ist $V C\varepsilon^2/2$,
der Kopplungsterm $-V\varepsilon eE$ und der elektrische Term $-V\kappa E^2/2$.
Eine analytische Randintegration der konstanten Traktion und des D-Flusses
prüft die nodalen Residuen unabhängig von der Matrixassembly.

\section{Handrechnung: geklemmter thermisch angeregter Grenzfall}
$C_{xx}=10$ Pa, $e_{xx}=2$ C/m$^2$, $\kappa_{xx}=4$ F/m,
$\alpha_x=0{,}01$/K, $T_{\rm ref}=300$ K, $T=310$ K.
Alle mechanischen DOFs sind fest, also $\varepsilon=0$ und
$\varepsilon_{{\rm th},x}=0{,}1$.
Bei $\chi=1$ und Ladungsfreiheit folgt
$0=2(-0{,}1)+4E_x$, somit $E_x=0{,}05$ V/m und bei
$\phi(0)=0$ die freie Ecke $\phi(1)=-0{,}05$ V.
Die Spannung ist $\sigma_x=-1{,}1$ Pa.
Bei $\chi=0$ ist $E_x=0$, $\phi=0$ und $\sigma_x=-1$ Pa.
Damit ist der elektrische Modellschalter beobachtbar und kein kosmetisches
Postprocessing. Der Thermopiezo-Test prüft das Potential unabhängig bis
$10^{-12}$ V.

\section{d33-Stab: zwei Anregungen und Leerlauf}\label{sec:d33}
Ein 1D-invarianter Stab in Polarisationsrichtung $z$ besitzt
$L=0{,}02$ m, $A=10^{-4}$ m$^2$, $C=10^9$ Pa, $e=5$ C/m$^2$ und
$\kappa=10^{-8}$ F/m. Die synthetischen diagonalen Tensoren und vorgegebenen
Querverschiebungen erlauben genau diese Reduktion:
$\sigma=Cu'+e\phi'$, $D=eu'-\kappa\phi'$, $\sigma'=D'=0$.
Für freie Mechanik und $\phi(L)=V$ gilt
\begin{equation}\label{eq:d33-voltage}
u(L)=-\frac eC V,\quad
Q_{L,\rm top}=\frac AL\left(\kappa+\frac{e^2}{C}\right)V,\quad
\mathcal H=-\tfrac12Q_{L,\rm top}V.
\end{equation}
Bei $V=100$ V erhält man $u(L)=-5\cdot10^{-7}$ m,
$Q_{L,\rm top}=1{,}75\cdot10^{-8}$ C, $\mathcal H=-8{,}75\cdot10^{-7}$ J.
Kurzschluss und Kraft $F=10$ N ergeben
$u(L)=FL/(AC)=2\cdot10^{-6}$ m und
$Q_{L,\rm top}=-eF/C=-5\cdot10^{-8}$ C.
Damit ist $u(L)/V=Q_{L,\rm top}/F=-e/C$, ein reziproker Cross-Compliance-Test.

Im statischen 1D-Leerlauf $D=0$ ist $C_D=C+e^2/\kappa=3{,}5\cdot10^9$ Pa.
Es folgen $u(L)=FL/(AC_D)=4\cdot10^{-6}/7$ m und
$\phi(L)=eFL/(A\kappa C_D)=2000/7$ V.
Nur in diesem invarianten Grenzfall ist eine punktweise Nullflussfläche von
selbst equipotentiell. Allgemeine Leiter benötigen die Kondensation.

\subsection*{Implementierung und unabhängige Nachweise.}
{\small\raggedright
\path{src/fem_toolbox/elements/piezoelectric.py},
\path{src/fem_toolbox/solvers/piezoelectric.py},
\path{src/fem_toolbox/solvers/thermopiezoelectric.py};
\path{tests/patch/test_piezo_element_patch.py},
\path{tests/patch/test_piezo_solver_patch.py},
\path{tests/unit/test_thermopiezoelectric.py},
\path{tests/regression/test_piezo_d33_benchmark.py} und die eigenständige
1D-Referenz \path{tests/reference/piezo_bar_1d.py}.
Diese Referenz importiert keine Produktionsoperatoren. Vier-Netze-Nachweise
für Kopplungen und Tet10 sind in Kapitel~\ref{chap:evidence} zugeordnet.

\par}
